\documentclass[DIN, pagenumber=false, fontsize=11pt, parskip=half]{scrartcl}

\usepackage[utf8]{inputenc}
\usepackage[T1]{fontenc}
\usepackage{template}
\usepackage{amsmath}

\DeclareMathOperator{\lcm}{lcm}

\titlehead{\centering\small
    Worked solutions to selected exercises from\\
    Richard Hammack, \emph{Book of Proof, edition 2.2}\\
\rule{0.4\linewidth}{0.4pt}}

\title{}
\author{} % Oscar A. Carrera
\date{}

\begin{document}
\maketitle

\setcounter{section}{7}

\section{Proofs Involving Sets}

\textbf{Use the methods introduced in this chapter to prove the following statements.}

\begin{problems}
    \problem Prove that $\set{6n : n \in \mathbb{Z}} = \set{2n : n \in \mathbb{Z}} \cap \set{3n : n \in \mathbb{Z}}$.
    \begin{proof}
        First we show $\set{6n : n \in \mathbb{Z}} \subseteq \set{2n : n \in \mathbb{Z}} \cap \set{3n : n \in \mathbb{Z}}$.
        Suppose $a \in \set{6n : n \in \mathbb{Z}}$.
        This means $a = 6n$ for some integer $n$.
        Therefore $a = 2(3n)$ and $a = 3(2n)$.
        From $a = 2(3n)$, it follows that $a$ is a multiple of 2,
        so $a \in \set{2n : n \in \mathbb{Z}}$.
        From $a = 3(2n)$, it follows that $a$ is a multiple of 3,
        so $a \in \set{3n : n \in \mathbb{Z}}$.
        Thus by definition of intersection of two sets,
        we have $a \in \set{2n : n \in \mathbb{Z}} \cap \set{3n : n \in \mathbb{Z}}$.
        Therefore $\set{6n : n \in \mathbb{Z}} \subseteq \set{2n : n \in \mathbb{Z}} \cap \set{3n : n \in \mathbb{Z}}$.

        Next we show $\set{2n : n \in \mathbb{Z}} \cap \set{3n : n \in \mathbb{Z}} \subseteq \set{6n : n \in \mathbb{Z}}$.\\ 
        Suppose $a \in \set{2n : n \in \mathbb{Z}} \cap \set{3n : n \in \mathbb{Z}}$.
        By definition of intersection of two sets, this means $a \in \set{2n : n \in \mathbb{Z}}$ and $a \in \set{3n : n \in \mathbb{Z}}$. 
        Since $a \in \set{2n : n \in \mathbb{Z}}$, we know $a = 2k$ for some $k \in \mathbb{Z}$.
        Thus $a$ is even.
        Since $a \in \set{3n : n \in \mathbb{Z}}$, we know $a = 3l$ for some $l \in \mathbb{Z}$.
        Since 3 is odd, if $l$ were odd then $3l$ would be odd; as $a = 3l$ is even, it follows that $l$ is even.
        Then $l = 2j$ for some $j \in \mathbb{Z}$.
        So we have $a = 3l = 3(2j) = 6j$.
        Since $j \in \mathbb{Z}$, it follows that $a \in \set{6n : n \in \mathbb{Z}}$.
        Thus $\set{2n : n \in \mathbb{Z}} \cap \set{3n : n \in \mathbb{Z}} \subseteq \set{6n : n \in \mathbb{Z}}$.

        Therefore, since  $\set{6n : n \in \mathbb{Z}} \subseteq \set{2n : n \in \mathbb{Z}} \cap \set{3n : n \in \mathbb{Z}}$ and $\set{2n : n \in \mathbb{Z}} \cap \set{3n : n \in \mathbb{Z}} \subseteq \set{6n : n \in \mathbb{Z}}$, it follows that $\set{6n : n \in \mathbb{Z}} = \set{2n : n \in \mathbb{Z}} \cap \set{3n : n \in \mathbb{Z}}$.
    \end{proof}

    \problem If $m,n \in \mathbb{Z}$, then $\set{x \in \mathbb{Z} : mn \mid x} \subseteq \set{x \in \mathbb{Z} : m \mid x} \cap \set{x \in \mathbb{Z} : n \mid x}$.
    \begin{proof}
        Suppose $a \in \set{x \in \mathbb{Z} : mn \mid x}$.
        Then $a = mnk$ for some $k \in \mathbb{Z}$.
        So we have $a = m(nk)$ and $a = n(mk)$.
        By definition of divisibility, $m \mid a$ and $n \mid a$.
        Thus $a \in \set{x \in \mathbb{Z} : m \mid x}$ and $a \in \set{x \in \mathbb{Z} : n \mid x}$.
        By definition of intersection of two sets, we have $a \in \set{x \in \mathbb{Z} : m \mid x} \cap \set{x \in \mathbb{Z} : n \mid x}$.
        Therefore $\set{x \in \mathbb{Z} : mn \mid x} \subseteq \set{x \in \mathbb{Z} : m \mid x} \cap \set{x \in \mathbb{Z} : n \mid x}$.
    \end{proof}

    \problem Suppose $A,B$ and $C$ are sets. Prove that if $A \subseteq B$, then $A - C \subseteq B - C$.
    \begin{proof}
        We prove the contrapositive: if $A - C \not\subseteq B - C$, then $A \not\subseteq B$.
        Assume $A - C \not\subseteq B - C$.
        Then there exists an element $x \in A - C$ such that $x \not\in B - C$.
        Since $x \in A - C$, we have $x \in A$ and $x \not\in C$.
        Since $x \not\in B - C$, it follows that $x \not\in B$ or $x \in C$.
        Because $x \not\in C$, it follows $x \not\in B$.
        Since $x \in A$ and $x \not\in B$, by definition of a subset, we have $A \not\subseteq B$.
    \end{proof}

    \problem If $A,B$ and $C$ are sets, then $A \cup (B \cap C) = (A \cup B) \cap (A \cup C)$.
    \begin{proof}
        Observe the following sequence of equalities.
        \begin{align*}
            A \cup (B \cap C) & = \set{x : x \in A \lor x \in B \cap C} && \text{(definition of $\cup$)}\\
                        & = \set{x : x \in A \lor (x \in B \land x \in C)} && \text{(definition of $\cap$)}\\
                        & = \set{x : (x \in A \lor x \in B) \land (x \in A \lor x \in C)} && \text{(distributive law)}\\
                        & = \set{x : x \in A \cup B \land x \in A \cup C} && \text{(definition of $\cup$)}\\
                        & = \set{x : x \in (A \cup B) \cap (A \cup C)} && \text{(definition of $\cap$)}\\
                        & = (A \cup B) \cap (A \cup C) && \text{(definition of set-builder)}
        \end{align*}
    \end{proof}

    \problem If $A$ and $B$ are sets in a universal set $U$, then $\overline{A \cap B} = \overline{A} \cup \overline{B}$.
    \begin{proof}
        Observe the following sequence of equalities.
        \begin{align*}
            \overline{A \cap B} & = \set{x \in U : x \notin A \cap B} && \text{(definition of complement)}\\
                                & = \set{x \in U : \neg(x \in A \land x \in B)} && \text{(definition of $\cap$)}\\
                                & = \set{x \in U : x \notin A \lor x \notin B} && \text{(DeMorgan's law)}\\
                                & = \set{x \in U : x \in \overline{A} \lor x \in \overline{B}} && \text{(definition of complement)}\\
                                & = \set{x \in U : x \in \overline{A} \cup \overline{B}} && \text{(definition of $\cup$)}\\
                                & = \overline{A} \cup \overline{B} && \text{(definition of set-builder)}
        \end{align*}
    \end{proof}

    \problem If $A,B$ and $C$ are sets, then $A - (B \cap C) = (A - B) \cup (A - C)$.
    \begin{proof}
        Observe the following sequence of equalities.
        \begin{align*}
            A - (B \cap C) & = \set{x : (x \in A) \land (x \notin B \cap C)} && \text{(definition of difference)}\\
                           & = \set{x : (x \in A) \land \lnot(x \in B \cap C)} && \text{(notation)}\\
                           & = \set{x : (x \in A) \land \lnot (x \in B \land x \in C)} && \text{(definition of $\cap$)}\\
                           & = \set{x : (x \in A) \land (\lnot(x \in B) \lor \lnot(x \in C))} && \text{(DeMorgan's law)}\\
                           & = \set{x : (x \in A) \land (x \notin B \lor x \notin C)} && \text{(notation)}\\
                           & = \set{x : (x \in A \land x \notin B) \lor (x \in A \land x \notin C)} && \text{(distributive law)}\\
                           & = \set{x : x \in A - B \lor x \in A - C} && \text{(definition of difference)}\\
                           & = \set{x : x \in (A - B) \cup (A - C)} && \text{(definition of $\cup$)}\\
                           & = (A - B) \cup (A - C) && \text{(definition of set-builder)}
        \end{align*}
    \end{proof}

    \problem If $A,B$ and $C$ are sets, then $(A \cup B) - C = (A - C) \cup (B - C)$. 
    \begin{proof}
        Observe the following sequence of equalities.
        \begin{align*}
            (A \cup B) - C & = \set{x : x \in (A \cup B) \land (x \notin C)} && \text{(definition of difference)}\\ 
                           & = \set{x : (x \in A \lor x \in B) \land (x \notin C)} && \text{(definition of $\cup$)}\\ 
                           & = \set{x : (x \in A \land x \notin C) \lor (x \in B \land x \notin C)} && \text{(distributive law)}\\ 
                           & = \set{x : x \in A - C \lor x \in B - C} && \text{(definition of difference)}\\
                           & = \set{x : x \in (A - C) \cup (B - C)} && \text{(definition of $\cup$)}\\
                           & = (A - C) \cup (B - C) && \text{(definition of set-builder)}
        \end{align*}
    \end{proof}

    \problem If $A,B$ and $C$ are sets, then $A \times (B \cup C) = (A \times B) \cup (A \times C)$.
    \begin{proof}
        Observe the following sequence of equalities.
        \begin{align*}
            A \times (B \cup C) & = \set{(x,y) : x \in A \land y \in B \cup C} && \text{(definition of $\times$)}\\
                                & = \set{(x,y) : x \in A \land (y \in B \lor y \in C)} && \text{(definition of $\cup$)}\\ 
                                & = \set{(x,y) : (x \in A \land y \in B) \lor (x \in A \land y \in C)} && \text{(distributive law)}\\ 
                                & = \set{(x,y) : (x,y) \in A \times B \lor (x,y) \in A \times C} && \text{(definition of $\times$)}\\
                                & = \set{(x,y) : (x,y) \in (A \times B) \cup (A \times C)} && \text{(definition of $\cup$)}\\ 
                                & = (A \times B) \cup (A \times C) && \text{(definition of set-builder)}
        \end{align*}
    \end{proof}

    \problem If $A,B$ and $C$ are sets, then $A \times (B - C) = (A \times B) - (A \times C)$.
    \begin{proof}
        Observe the following sequence of equalities.
        \begin{align*}
            A \times (B - C) & = \set{(x,y) : x \in A \land y \in B-C} && \text{(definition of $\times$)}\\ 
                             & = \set{(x,y) : x \in A \land (y \in B \land y \notin C)} && \text{(definition of difference)}\\
                             & = \{(x,y) : (x \in A \land y \in B \land x \notin A) && \text{(first disjunct is false)}\\
                             & \hphantom{{}= \{(x,y) :{}} \lor (x \in A \land y \in B \land y \notin C)\}\\
                             & = \set{(x,y) : (x \in A \land y \in B) \land (x \notin A \lor y \notin C)} && \text{(distributive law)}\\
                             & = \set{(x,y) : x \in A \land y \in B \land \lnot(x \in A \land y \in C)} && \text{(DeMorgan's law)}\\
                             & = \set{(x,y) : (x,y) \in A \times B \land (x,y) \notin A \times C} && \text{(definition of $\times$)}\\
                             & = \set{(x,y) : (x,y) \in (A \times B) - (A \times C)} && \text{(definition of difference)}\\
                             & = (A \times B) - (A \times C) && \text{(definition of set-builder)}
        \end{align*}
    \end{proof}

    \problem Prove that $\set{9^n : n \in \mathbb{Q}} = \set{3^n : n \in \mathbb{Q}}$. 
    \begin{proof}
        First we prove $\set{9^n : n \in \mathbb{Q}} \subseteq \set{3^n : n \in \mathbb{Q}}$.
        Assume $a \in \set{9^n : n \in \mathbb{Q}}$.
        This means $a = 9^n$ for some $n \in \mathbb{Q}$.
        Thus $a = 9^n = (3^2)^n = 3^{2n}$.
        Since $2n \in \mathbb{Q}$, we have $a \in \set{3^n : n \in \mathbb{Q}}$.
        Therefore $\set{9^n : n \in \mathbb{Q}} \subseteq \set{3^n : n \in \mathbb{Q}}$.

        Now we prove $\set{3^n : n \in \mathbb{Q}} \subseteq \set{9^n : n \in \mathbb{Q}}$.
        Assume $a \in \set{3^n : n \in \mathbb{Q}}$.
        This means $a = 3^n$ for some $n \in \mathbb{Q}$.
        Thus $a = 3^n = (9^{1/2})^n = 9^{n/2}$.
        Since $n/2 \in \mathbb{Q}$, we have $a \in \set{9^n : n \in \mathbb{Q}}$.
        Therefore $\set{3^n : n \in \mathbb{Q}} \subseteq \set{9^n : n \in \mathbb{Q}}$.

        Therefore, since $\set{9^n : n \in \mathbb{Q}} \subseteq \set{3^n : n \in \mathbb{Q}}$ and $\set{3^n : n \in \mathbb{Q}} \subseteq \set{9^n : n \in \mathbb{Q}}$, it follows that $\set{9^n : n \in \mathbb{Q}} = \set{3^n : n \in \mathbb{Q}}$.
    \end{proof}

    \problem Let $A$ and $B$ be sets. Prove that $A \subseteq B$ if and only if $A \cap B = A$.
    \begin{proof}
        First we prove if $A \subseteq B$, then $A \cap B = A$.\\ 
        Suppose $A \subseteq B$. Then every element of $A$ is also an element of $B$. Thus an element is in both $A$ and $B$ exactly when is in $A$.
        Therefore $A \cap B = A$.

        Conversely, we prove if $A \cap B = A$, then $A \subseteq B$.\\ 
        Suppose $A \cap B = A$ and let $x \in A$.
        Since $A = A \cap B$, it follows that $x \in B$.
        Thus every element of $A$ is also an element of $B$.
        Therefore $A \subseteq B$.
    \end{proof} 


\end{problems}

\end{document}
